\chapter{BIBLIOGRAFÍA Y ABREVIATURAS}
\section{\emph{Historia, geografía, etc.}}

\subsection{} \begin{description} 
\item[CABRERA, Ángel] \textit{Mamíferos}, Madrid, 1914.
\item[CARO BAROJA, J.] \textit{Pueblos del norte de la Península Ibérica}, Madrid, 1943.
\item[COLMEIRO, Miguel] \textit{Enumeraciones y revisión de las plantas de la Península Hispano-Lusitana e Islas Baleares, etc.}, Madrid, 1885-89.
\item[Cossío, J. M.ª de, y MAZA SOLANO, T.] \textit{Romancero popular de la Montaña}, Santander, t. I, 1933, t. II, 1934. 
\item[ESCAGEDO SALMÓN, M.] \textit{Costumbres pastoriles cántabro-montañesas}, Santander, 1921.
\item[ESPERÓN, Antolín] `El Pasiego', \textit{Seminario Pintoresco Español}, 1851, 390-392.
\item[ESPINOSA, A.M.] \textit{Cuentos populares españoles}, Stanford University, 3 t., 1923, 1924, 1926. 
\item[GARCÍA BELLIDO, A.] \textit{Cantabria Romana}, Santander, 1952. 
\item[GARCÍA BELLIDO, A.] \textit{España y los españoles hace dos mil años según la geografía de Strabón}, Madrid, 1945.
\item[GARCÍA LOMAS, G. A.] \textit{Los Pasiegos. Estudio etnográfico e histórico}, Santander, 1960.
\item[G(IL), E(nrique)] \textit{Semanario Pintoresco Español}, 1839, 201-203. 
\item[GUINEA LOPEZ, E.] \textit{Geografía botánica de Santander}, Santander, 1953. 
\item[JUSUÉ, E.] ed., \textit{Cartulario de Santillana del Mar}, 1912. 
\item[LASAGA LARRETA, G.] \textit{Compilación histórica, biográfica y marítima de la provincia de Santander}, Cádiz, 1865.
\item[MADOZ, P.] \textit{Diccionario Geográfico de España, etc.}, Madrid, 1846.
\item[NOMENCLÁTOR] de la provincia de Santander, 1960.
\item[ORTIZ DE LA TORRE, Elías] `La etnografía en la obra de Pereda', \textit{Altamira}, I. 1934.
\item[DE LOS RÍOS, Amador] \textit{España, sus monumentos y artes, su naturaleza e historia}, t. \textit{Santander}, Barcelona, 1891. 
\item[SÁNCHEZ BELDA, L.] \textit{Cartulario de Sto. Toribio de Liébana}, ed. y estudio de $\sim$, 1948.
\item[SCHULTEN, A.] \textit{Los cántabros y astures y su guerra con Roma}, Madrid, 1943. 
\item[TERÁN, Manuel de] `Vaqueros y cabañas en los Montes de Pas', \textit{Estudios Geográficos}, VIII, 1947.
\item[TOVAR, A.] \textit{Cantabria romana. El problema de los cántabros y su origen o lo que la lingüística nos enseña sobre los antiguos Cántabros}, Publicaciones de la Universidad Internacional Menéndez Pelayo, Madrid, 1955.
\end{description}

\section{\emph{Léxico, etimología, etc.}}

\subsection{} \begin{description}
\item[ACEVEDO Y HUELVES, I, y FERNÁNDEZ, M.] \textit{Vocabulario del bable de occidente}, Madrid, 1932.
\item[ALCALÁ VENCESLADA, A.] \textit{Vocabulario Andaluz}, Andújar, 1933.
\item[Alcalde del Río, 'Contribución' = ALCALDE DEL RÍO, H.] `Contribución al léxico montañés', \textit{Revista de Santander}, V, 1932, 238.
\item[Baráibar, Álava = BARÁIBAR Y ZUMÁRRAGA, F.] \textit{Vocabulario de palabras usadas en Álava}, Madrid, 1903.
\item[Borao = BORAO, Jerónimo] \textit{Diccionario de voces aragonesas}, 2.ª ed., Zaragoza, 1908.
\item[Calderón Escalada, I = CALDERÓN ESCALADA, J.] `Voces, en su mayor parte nombres de cosas, de uso corriente en estos valles altos de la provincia de Santander, que no están recogidas en el Diccionario de la Lengua Española', \textit{BRAE}, XXIII, 1953, 295-304. 
\item[Calderón Escalada, II = CALDERÓN ESCALADA, J.] `Voces, en su mayor parte nombres de cosas, de uso corriente en estos valles altos de la provincia de Santander, que no están recogidas en el Diccionario de la Lengua Española', \textit{BRAE}, XXV, 1946, 379-397.
\item[Caminals, `Guadilla de Villamar' = FUENTE CAMINALS, J. de la] `Algunas palabras de Guadilla de Villamar (Burgos)', \textit{RDTrP}, V, 1949, 146-151.
\item[Campuzano, `Apuntes' = GONZÁLEZ CAMPUZANO, J.] 'Apuntes para un vocabulario montañés', \textit{BBMP}, II, 1920. 
\item[CASTRO, A.] `Adiciones hispánicas al Diccionario etimológico de W. Meyer-Lübke', \textit{RFE}, V, 1918, 21-42.
\item[CÉSAR MORÁN, P.] `Vocabulario del Concejo de Lomba (León)', \textit{BRAE}, XXX, 1950, 155-168, 313-330 y 439-456.
\item[Coll y Altabás =  COLL Y ALTABÁS, B.] `Colección de voces usadas en la Litera', apéndice al diccionario de Borao, cf. supra. 
\item[Corominas = COROMINAS, J.] \textit{Diccionario crítico etimológico de la lengua castellana}, Berna, 1954-1957.
\item[CRUZ MARTÍNEZ, J. de la] 'Algunas voces de Hornillayuso (Burgos)',  \textit{RDTrP}, XVI, 1961, 174-180.
\item[DRAE = \textit{Diccionario de la Real Academia Española}] 18.ª ed., 1956.
\item[G. de Diego, \textit{Contribución} = GARCÍA DE DIEGO, V.] \textit{Contribución al diccionario hispánico etimológico}, Madrid, 1943. 
\item[García Lomas = GARCÍA LOMAS Y GARCÍA LOMAS, G.A.] \textit{El lenguaje popular de las montañas de Santander}, Santander, 1949.
\item[GARCÍA LOMAS Y GARCÍA LOMAS, G.A.] \textit{Estudio del dialecto popular montañés}, San Sebastián, 1922.
\item[Goicoechea, \textit{Vocab. riojano} = GOICOECHEA, C.]  \textit{Vocabulario riojano}, Madrid, 1961.
\item[Gutiérrez Cuñado, `Tierra de Campos' = GUTIÉRREZ CUÑADO, A.] `El léxico de Tierra de Campos', \textit{BRAE}, XXX, 1950, 257-262.
\item[Huidobro, \textit{Palabras} = HUIDOBRO, E.] `Palabras, giros y bellezas del lenguaje popular de la Montaña', Santander, 1907.   
\item[IRIBARREN, J. M.] \textit{Vocabulario navarro}, Pamplona, 1952.  
\item[KRUGER, F.] \textit{El léxico rural del noroeste peninsular}, Madrid, 1947. 
\item[Lamano = LAMANO Y BENEITE, J. de] \textit{El dialecto vulgar salmantino}, Salamanca, 1915.   
\item[López Barrera, \textit{Cuenca} = LÓPEZ BARRERA, J.] \textit{Arcaísmos y barbarismos en la provincia de Cuenca. Estudios de semántica regional}, Cuenca, 1912. 
\item[MANRIQUE, G.] `Vocabulario popular comparado de los valles de Duero y del Ebro', \textit{RDTrP}, XII, 1956, 3-53. 
\item[MORAL MORAL, M.] `Voces de Gumiel de Hizán (Burgos)', \textit{RDTrP}, XVII, 1961, 374-383.
\item[MORALES, G.] \textit{La Montaña, Añoranza}, Madrid, 1919 (contiene un corto vocabulario de modismos montañeses).  
\item[PARDO ASSÓ, J.] \textit{Nuevo diccionario aragonés}, Zaragoza, 1938. 
\item[Rato = RATO Y HEVIA, A.] \textit{Vocabulario de las palabras y frases bables que se hablaron antiguamente y de las que hoy se hablan en el principado de Asturias.} Compendio gramatical. Madrid, 1891. 
\item[RENERO, V.] `Formas dialectales y toponímicas de Cantabria', \textit{Altamira}, 1947.  
\item[\textit{REW} = MEYER-LÜBKE, W.] \textit{Romanisches Etimologisches Wörterbuch}, Heidelberg, 1924. 
\item[ROHLFS, G.] \textit{La diferencia léxica de las lenguas románicas} (tr. de M. Alvar), Madrid, 1960. 
\item[SIERRA PANDO, J.] `Vocabulario santanderino', BBMP, III, 1921.  
\item[Vergara, B. y S. = VERGARA MARTÍN, G. M.] \textit{Materiales para la formación de un vocabulario de palabras empleadas en Burgos y Santander}, publicado como apéndice de la obra siguiente: 
\item[Vergara, \textit{Segovia} = VERGARA MARTÍN, G. M.] \textit{Materiales para la formación de un vocabulario de palabras usadas en Segovia}, Madrid, 1921.
\item[Vigón = VIGÓN, B.] \textit{Vocabulario dialectológico del concejo de Colunga}, 2.ª ed., Madrid, 1955.
\end{description}


\section{\emph{Monografías dialectológicas.}}

\subsection{} \begin{description}
\item[Alonso Garrote, \textit{Maragatería} = ALONSO GARROTE, S.] \textit{El dialecto vulgar leonés hablado en la Maragatería y la tierra de Astorga}, 2.ª ed., Madrid, 1947.
\item[Alvar, \textit{Campo de Jaca} = ALVAR, M.] \textit{El habla del Campo de Jaca}, Salamanca, 1947.   
\item[Alvarez, \textit{Babia y Laciana} = ALVAREZ, G.] \textit{El habla de Babia y Laciana}, Madrid, 1949. 
\item[Alvarez Fernández Cañedo, \textit{Cabrales} = ALVAREZ FERNÁNDEZ CAÑEDO, J.] \textit{El habla y la cultura popular de Cabrales (Este de Asturias)}, Madrid, 1963.
\item[Casado Lobato, \textit{Cabrera Alta} = CASADO LOBATO, M.C.] \textit{El habla de la Cabrera Alta}, Madrid, 1948. 
\item[DÍAZ CASTAÑÓN, C.] \textit{El bable de `El cabo Peñas'}, Oviedo, 1966.
\item[ESPINOSA, A. M.] \textit{Estudios sobre el español de Nuevo Méjico}, 2 tomos, anotados por A. Castro y A. Rosenblat, Buenos Aires, 1930 y 1946. 
\item[FERNÁNDEZ, J. A.] \textit{El habla de Sisterna}, Madrid, 1960. 
\item[FERNÁNDEZ GONZÁLEZ, A. R.] \textit{El habla y la cultura popular de Oseja de Sajambre}, Oviedo, 1959.  
\item[G. de Diego, `Dialectalismos' = GARCÍA DE DIEGO, V.] `Dialectalismos castellanos', \textit{RFE}, III, 1916, 300-318. 
\item[GARCÍA DE DIEGO, V.] `El castellano como ejemplo dialectal y sus dialectos internos', \textit{RFE}, XXXIV, 1950, 107-124. 
\item[KRÜGER, F.] `Mezcla de dialectos', \textit{Homenaje a Menéndez Pidal}, t. II, Madrid, 1925, 121-166. 
\item[KRÜGER, \textit{San Ciprián} = KRÜGER, F.] \textit{El dialecto de San Ciprián de Sanabria}, Madrid, 1923. 
\item[Llorente Maldonado, \textit{La Ribera} = LLORENTE MALDONADO DE GUEVARA, A.] \textit{Estudio sobre el habla de la Ribera}, Salamanca, 1947.
\item[Menéndez Pidal, \textit{Lena} = MENÉNDEZ PIDAL, R.] \textit{Notas acerca del bable de Lena}, publicado en nueva edición, junto con \textit{El dialecto leonés} (cf. abajo), por el IDEA, Oviedo, 1962.  
\item[Monge, `Puebla de Hijar' = MONGE, F.] `El habla de la Puebla de Hijar (Aragón)', \textit{RDTrP}, VII, 1951, 187-241.
\item[Neira, \textit{Lena} = NEIRA MARTÍNEZ, J.] \textit{El habla de Lena}, Oviedo, 1955.
\item[Ollé, \textit{Bureba} = GONZÁLEZ OLLÉ, F.] \textit{El habla de la Bureba}, Madrid, 1964.
\item[Ollé, `Burgos' = GONZÁLEZ OLLÉ, F.] `El habla de Burgos como modelo idiomático en la historia de la lengua española y su situación actual', \textit{Presente y futuro de la lengua española}, Madrid, 1964, 227-237.
\item[Ollé, `Mena' = GONZÁLEZ OLLÉ, F.] `Características fonéticas y léxico del valle de Mena (Burgos)', \textit{BRAE}, XL, 1960, 67-85.
\item[Ollé, `Quintanillabón' = GONZÁLEZ OLLÉ, F.] `El habla de Quintanillabón (Burgos)', \textit{RDTrP}, IX, 1953, 3-65.
\item[PENNY, R.J.] `Mass-nouns and metaphony in the dialects of Northwestern Spain', \textit{Archivum Linguisticum} (en prensa).
\item[PENNY, R.J.] `Vowel-harmony in the speech of Montes de Pas (Santander)', \textit{Orbis}, XVIII, 1969, 148-166. 
\item[Rodríguez Castellano, \textit{Alto Aller} = RODRÍGUEZ CASTELLANO, L.] \textit{La variedad dialectal del Alto Aller}, Oviedo, 1952. 
\item[RODRÍGUEZ CASTELLANO, L.] \textit{Aspectos del bable de occidente}, Oviedo, 1954. 
\item[Rohlfs, \textit{Le Gascon} = ROHLFS, G.] \textit{Le Gascon. Études de philologie pyrénéenne}, \textit{ZfRPh}, Beiheft 85, Halle, 1935.
\item[Sevilla, `Cespedosa' = SÁNCHEZ SEVILLA, P.] `El habla de Cespedosa de Tormes'. \textit{RFE}, XV, 1928, 131-172 y 244-288.
\item[Torre, `Cuéllar' = TORRE, Alfonso de la] `El habla de Cuéllar (Segovia)', \textit{BRAE}, 1951, 133-164 y 501-513. 
\item[Zamora, \textit{Mérida} = ZAMORA VICENTE, A.] \textit{El habla de Mérida y sus cercanías}, Madrid, 1943.
\end{description}


\section{\emph{Otros trabajos de dialectología.}}
\setcounter{subsection}{18}
\subsection{a} \begin{description}
\item[ALONSO, A.] sobre Post, Anita C., \textit{Southern Arizona Spanish Phonology}, \textit{RFE}, XXII, 1935, 67-72.
\item[ALONSO, D.] `Metafonía y neutro de materia en España (sobre un fondo italiano)', \textit{ZfRPh}, LXXIV, 1958, 1-24.
\item[ALONSO, D.]`Metafonía, neutro de materia y colonización suditaliana en la península hispánica', \textit{Enciclopedia lingüística hispánica}, t. I (Suplemento), 1962, 105-154.
\item[ALVAR, M.] \textit{El dialecto aragonés}, Madrid, 1953. 
\item[BLAYBOCK, C.] `Hispanic Metaphony',  \textit{Romance Philology}, XVIII, 1965, 253-271.
\item[CATALÁN MENÉNDEZ PIDAL, D. y GALMÉS DE FUENTES, A.] `Un límite lingüístico', \textit{RDTrP}, I, 1946. 
\item[CATALÁN MENÉNDEZ PIDAL, D. y GALMÉS DE FUENTES, A.] `La diptongación en leonés', \textit{Archivum}, IV, 1954, 87-147.
\item[CATALÁN MENÉNDEZ PIDAL, D. y GALMÉS DE FUENTES, A.] `Inflexión de las vocales tónicas junto al cabo de Peñas', \textit{RDTrP}, IX, 1953.
\item[ESPINOSA, A. M., y RODRÍGUEZ CASTELLANO, L.] `La aspiración de la h en el sur y oeste de España', \textit{RFE}, 1936, 225-254 y 337-378. 
\item[GALMÉS DE FUENTES, A.] `Más datos sobre la inflexión metafonética en el centro-sur de Asturias', \textit{Trabajos sobre el dominio románico leonés}, t. II, 13-25.
\item[GARCÍA ALVAREZ, M. T. C.], `La inflexión vocálica en el bable de Bimenes', \textit{BIDEA}, XXIV, 1955, 123-146.
\item[GARCÍA ALVAREZ, M. T. C.], `Morfología verbal en el bable de Bimenes', \textit{Archivum}, X, 1960, 405-424.
\item[GARCÍA DE DIEGO, V.] `Encuestas dialectales', \textit{RDTrP}, VII, 1951, 3-16. 
\item[GARCÍA DE DIEGO, V.] \textit{Manual de dialectología española}, Madrid, 1959.
\item[GARVENES, F.], `La metafonía en Cabrales (oriente de Asturias)', \textit{BIDEA}, 1960, XIV, 241-244.
\item[GRANDA GUTIÉRREZ, G. de], `Observaciones sobre el sistema morfológico del nombre en asturiano', \textit{RFE}, XLVI, 1963, 97-120. 
\item[GRANDA GUTIÉRREZ, G. de], `Las vocales finales del dialecto leonés', \textit{Trabajos sobre el dominio románico leonés}, t. II, Madrid, 1960, 20-117.
\item[HANSSEN, F.], \textit{Estudio sobre  la conjugación leonesa}, \textit{Archivum}, X, 1960, 271-318.
\item[INDURAIN, F.] \textit{Contribución al estudio del dialecto navarro-aragonés antiguo}, Zaragoza, 1945.
\item[LÓPEZ SANTOS, L.], `La diptongación en leonés', \textit{Archivum}, X, 1960, 271-318.
\item[MALKIEL, Y.], sobre Rodríguez Castellano, L., \textit{La variedad dialectal del Alto Aller}, \textit{Language}, XXX, 1954, 126-153.
\item[MENÉNDEZ PIDAL, R.], `Colonización suditálica en España', \textit{BRAE}, XXXIV, 1954, 165-216.
\item[MENÉNDEZ PIDAL, R.], \textit{El dialecto leonés}, nueva edición, Oviedo, 1962.
\item[MENÉNDEZ PIDAL, R.], `Pasiegos y Vaqueiros. Dos cuestiones de geografía lingüística', \textit{Archivum}, IV, 1954, 7-44.
\item[MENÉNDEZ PIDAL, R. y TOVAR, A.],  `Los sufijos con -rr- en España y fuera de ella, especialmente en la toponimia', \textit{BRAE}, XXXVIII, 1958, 161-214. 
\item[Múgica, \textit{Dialectos} = MÚGICA, P. de], \textit{Dialectos castellanos: montañés, vizcaíno, aragonés}, Berlín, 1892.
\item[NAVARRO TOMÁS, T.], `Datos antiguos sobre pronunciación asturiana', \textit{RFE}, VII, 1920, 382-383.
\item[PAIVA BOLÉO, M.] sobre Zamora Vicente, A., \textit{El habla de Mérida y sus cercanías},  \textit{Revista Portuguesa de Filología}, I, 531-543. 
\item[RODRÍGUEZ CASTELLANO, L.], `El posesivo en el dialecto asturiano', \textit{BIDEA}, XXXI, 1957, 171-187.
\item[RODRÍGUEZ CASTELLANO, L.], `Más datos sobre la inflexión vocálica en la zona centro-sur de Asturias', \textit{BIDEA}, XXIV, 1955, 123-140.
\item[ROHLFS, G.] \textit{Estudios sobre geografía lingüística de Italia} (tr. y prólogo de M. Alvar), Granada, 1952.
\item[ROSENBLAT, A.], `Morfología del género en español. Comportamiento de las terminaciones -o, -a', \textit{NRPH}, XVI, 1962, 1-2.
\item[SCHÜRR, F.], `Über Umlaut und innere Flexion im Asturischen', \textit{Vox Romanica}, XVII, 1958, 260-266.
\item[SCHÜRR, F.], `La diptongación ibero-románica', \textit{RDTrP}, VII, 1951, 379 \textit{et seq.}
\item[SOJO y LOMBA, F. de], \textit{De re toponímica. Comunicaciones en Cantabria}. Publicaciones de la Real Sociedad Geográfica, serie B, núm. 190, 1947, 7-71.
\item[STAAF, E.], \textit{Etude sur l'ancien dialecte léonais}, upsala, 1907.
\item[TILANDER, G.], `La terminación -i por -e en los poemas de Berceo', \textit{RFE}, XXIV, 1940. \textit{Trabajos sobre el dominio románico leonés}, Madrid, t. I, 1957, t. II, 1960.
\item[ZAMORA VICENTE, A.], \textit{Dialectología española}, Madrid, 1960.
\end{description}

\section{\emph{Obras generales.}}

\setcounter{subsection}{18}
\subsection{b} \begin{description}
\item[ALONSO, A.], \textit{Estudios lingüísticos. Temas españoles}, Madrid, 1951.
\item[BADÍA MARGARIT, A.], `Ensayo de una sintaxis histórica de tiempos', \textit{BRAE}, XXVIII, 1948, 281-300 y 393-410, y XXIX, 1949, 15-29.
\item[BOURCIEZ, E.], \textit{Éléments de Linguistique Romane}, 2.ª ed., París, 1923. \textit{Enciclopedia lingüística hispánica}, Madrid, t. I (Suplemento), 1962.
\item[JUNGEMANN, F.], \textit{La teoría del sustrato y los dialectos hispano-romances y gascones}, Madrid, 1955.
\item[MALKIEL, Y.], `Filología española y lingüística general', \textit{Actas del primer congreso internacional de hispanistas}, Oxford, 1964, 107-126.
\item[MENÉNDEZ PIDAL, R.], \textit{Castilla. La tradición. El idioma}, 3.ª ed., Madrid, 1955.
\item[MENÉNDEZ PIDAL, R.], \textit{Documentos lingüísticos de España, I: Reino de Castilla}, Madrid, 1910.
\item[MENÉNDEZ PIDAL, R.], \textit{Manual de gramática histórica española}, Madrid, 10.ª ed., 1958.
\item[MENÉNDEZ PIDAL, R.], `Modo de obrar el substrato lingüístico', \textit{RFE},  XXXIV, 1950.
\item[Menéndez Pidal, \textit{Orígenes} = MENÉNDEZ PIDAL, R.], \textit{Orígenes del español},  5.ª ed., Madrid, 1964.
\item[MENÉNDEZ PIDAL, R.], \textit{Toponimia prerrománica hispana}, Madrid, 1952.
\item[MEYER-LÜBKE, W.], \textit{Grammaire des Langues Romanes} (trad. française par Eugene Rabiet), Nueva York, 1923.
\item[Navarro Tomás, T., \textit{Manual} = NAVARRO TOMÁS, T.], \textit{Manual de pronunciación española}, 10.ª ed., Madrid, 1961.
\item[NAVARRO TOMÁS, T.], `La metafonía vocálica', \textit{RFE}, X, 1923, 26-56.
\item[PEREDA, J. M. de], \textit{Obras Completas}, Madrid, 1933.
\end{description}

\tableofcontents

